\documentclass[11pt]{article}
\usepackage[margin=1in]{geometry}
\usepackage{amsmath,amssymb,amsthm,booktabs,tabularx}
\usepackage[T1]{fontenc}
\usepackage{lmodern,microtype}
\usepackage[colorlinks=true,allcolors=blue]{hyperref}
\emergencystretch=2em
\newtheorem{theorem}{Theorem}
\newtheorem{lemma}[theorem]{Lemma}
\newtheorem{corollary}[theorem]{Corollary}
\newtheorem{proposition}[theorem]{Proposition}
\newcommand{\rank}{\operatorname{rank}}
\newcommand{\im}{\operatorname{im}}
\title{Quantitative null tensor summands for\protect\\He--Tubbenhauer's Conjecture 4.2\protect\\
\large arXiv:2508.04054v1, p.~17:\protect\\extending their 2025 theorem}
\author{Haobo Ma\thanks{ChronoAI Pte Ltd, Singapore.}\and
Wenlin Zhang\thanks{National University of Singapore; ChronoAI Pte Ltd, Singapore.}}
\date{5 November 2026}
\begin{document}
\maketitle
\begin{abstract}
Extending the right-idempotent-generated case proved by He and Tubbenhauer
in 2025, we prove their Conjecture 4.2 in \emph{Tensor powers of
representations of (diagram) monoids}, arXiv:\allowbreak2508.04054v1, p.~17:
a finite-dimensional representation of a finite semigroup, with every element
acting singularly, has a one-dimensional zero-action direct summand in a positive
tensor power, and its indecomposable-summand count is asymptotic to the tensor
dimension. The proof works over every field. If the original action ranks are
\(r_s\) and the dimension is \(d\), the deficit in the summand count is at most
\(2\sum_s r_s^n\). A classical common-kernel/common-image splitting gives the exact
multiplicity of the null summand. We give nilpotent-chain formulas and an exact
recurrence for the published example. The same splitting proves the related monoid
Conjecture~3A.2 when the unit-group order is invertible in the field. A two-element
right-zero semigroup distinguishes the literal equality printed in
Proposition~4.4 from its asymptotic replacement.
\end{abstract}

\section{The question and a splitting lemma}
He and Tubbenhauer \cite[Conjecture 4.2, p.~17]{HT25} conjecture that
``\(Z\) is a summand of \(V^{\otimes m}\)'' for some \(m\geq1\), together
with the summand-count asymptotic below. We prove both assertions.
Their field is algebraically closed, their semigroup is finite, nonempty and
without identity, and their modules are finite dimensional and nonunital.
Write \(Z\) for the one-dimensional representation on which every element acts
as zero. In their notation, \(a(n)=d^n\) means \(b(n)/d^n\to1\), where
\(b(n)\) counts indecomposable summands with multiplicity; it does not mean
\(b(n)=d^n\) for each \(n\) \cite[\S1, p.~1]{HT25}.

He and Tubbenhauer proved the right-idempotent-generated case in December 2025
\cite[v2, Theorem 4A.14(e)--(f), p.~20; v3, p.~22]{HT26}, assuming algebraic
closure. Their related monoid theorem is \cite[Theorem 4A.12(c)--(d)]{HT26}.
The finite singular-semigroup argument here uses only ranks and removes both
right-idempotent-generation and algebraic-closure assumptions. It does not
extend their infinite-semigroup results beyond their stated hypotheses.

The splitting principle is classical: McAlister's null-summand criterion
\cite[Definition~1.4, pp.~195--196; Corollary~2.6, p.~201]{McA71}
and Juteau--Mautner--Williamson's multiplicity pairing
\cite[v1, \S3.2.1, p.~23, (3.1); v3, Lemma~3.1, pp.~22--23]{JMW}
give antecedents for the lemma below. Our application combines that principle
with tensor-rank decay. General semigroup representations already satisfy
\(b(n)^{1/n}\to d\) \cite[Theorem~1.4(a), Remark~1.5(a)]{COT}; the analogous
pointwise-quiver result is \cite[Theorem~A]{LZ}. A root limit alone does not give
the ratio limit proved here: \((d^n/\sqrt n)^{1/n}\to d\), whereas
\((d^n/\sqrt n)/d^n\to0\).

Let \(A=(A_s)_{s\in F}\) be a finite family of endomorphisms of a finite-dimensional
vector space \(W\) over any field \(k\). An invariant subspace means invariant
under every \(A_s\). Define
\[
 K=\bigcap_{s\in F}\ker A_s,\qquad
 I=\sum_{s\in F}\im A_s,\qquad R=\sum_{s\in F}\rank A_s.
\]
Let \(\mu(W)\) be the largest dimension of a zero-action invariant direct
summand of \(W\).

\begin{lemma}\label{lem:split}
There is an invariant decomposition \(W=L\oplus U\), with \(L\subseteq K\)
and \(I\subseteq U\), such that
\begin{equation}\label{eq:mu}
 \mu(W)=\dim L=\dim K-\dim(K\cap I)\geq\dim W-2R.
\end{equation}
The quantity \(\mu\) is additive on invariant direct sums and equals the number
of copies of \(Z\) in any indecomposable decomposition.
\end{lemma}
The multiplicity-pairing interpretation is explicit: regard \(W\) as a module
over the free unital algebra \(k\langle x_s:s\in F\rangle\), with \(x_s\)
acting by \(A_s\). Then
\(\operatorname{Hom}(Z,W)=K\) and
\(\operatorname{Hom}(W,Z)=(W/I)^*\).
Composition is evaluation, with left radical \(K\cap I\). Thus the identity
in \eqref{eq:mu} is a specialization of \cite{JMW}. We give its direct proof.
\begin{proof}
Choose \(L\) with \(K=(K\cap I)\oplus L\). Then \(L\cap I=0\), so a basis of
\(I\), together with a basis of \(L\), extends to a basis of \(W\).
Let \(U\) be the span of the basis of \(I\) and the additional basis vectors.
Thus \(W=L\oplus U\). Every operator kills \(L\), and every operator maps
\(U\) into \(I\subseteq U\). This gives the claimed invariant splitting.

Conversely, if \(W=L'\oplus U'\) is invariant and every operator kills \(L'\),
then \(L'\subseteq K\) and \(A_s(W)=A_s(U')\subseteq U'\). Hence
\(L'\cap I=0\), and the map \(L'\to K/(K\cap I)\) is injective. This proves
the equality in \eqref{eq:mu}. The map
\(W\to\bigoplus_{s\in F}\im A_s\), \(w\mapsto(A_sw)_s\), has kernel \(K\),
so \(\dim W-\dim K\leq R\). Also \(\dim I\leq R\). Substitution proves
the inequality.

For an invariant decomposition \(W=\bigoplus_jW_j\), each operator is block
diagonal. Consequently \(K=\bigoplus_jK_j\), \(I=\bigoplus_jI_j\), and
\(K\cap I=\bigoplus_j(K_j\cap I_j)\). The equality in \eqref{eq:mu} gives
additivity. If an indecomposable \(W_j\) has \(\mu(W_j)>0\), the constructed
splitting forces it to have zero action; indecomposability then forces its
dimension to be one. Conversely, a copy of \(Z\) has \(\mu=1\).
\end{proof}

\section{Tensor powers and quantitative consequences}
\begin{theorem}\label{thm:main}
Let \(S\) be a nonempty finite semigroup and
\(\rho:S\to\operatorname{End}_k(V)\) a multiplicative map, where \(V\) is finite
dimensional over any field. Suppose every \(\rho(s)\) is singular. Put
\(d=\dim V\), \(r_s=\rank\rho(s)\), and \(t=|S|\). Then \(d\geq1\).
For every \(n\geq1\), the diagonal action on \(V^{\otimes n}\) has a finite
indecomposable decomposition, and every such decomposition, with \(b(n)\)
factors, satisfies
\begin{equation}\label{eq:bounds}
 d^n-2\sum_{s\in S}r_s^n\ \leq\ \mu(V^{\otimes n})\ \leq\ b(n)\ \leq\ d^n.
\end{equation}
In particular, \(b(n)/d^n\to1\), and a copy of \(Z\) is a direct summand
whenever
\begin{equation}\label{eq:threshold}
 n\geq1,\qquad 2\sum_{s\in S}r_s^n<d^n.
\end{equation}
Such an exponent always exists.
\end{theorem}
\begin{proof}
The unique endomorphism of the zero space is invertible, so the singularity
assumption and \(S\ne\varnothing\) imply \(d\geq1\).
The diagonal operator associated to \(s\) is \(\rho(s)^{\otimes n}\).
To compute its rank, factor \(\rho(s)\) as a surjection onto its image followed
by the image inclusion. Tensoring these maps preserves surjectivity and
injectivity over a field, as can be seen by choosing bases and linear sections.
Its image is therefore the embedded space \((\im\rho(s))^{\otimes n}\), of
dimension \(r_s^n\). Also \(\dim V^{\otimes n}=d^n\).

Apply Lemma~\ref{lem:split} to this family of tensor operators. To obtain an
indecomposable decomposition of any invariant finite-dimensional space, use
induction on its dimension: the zero space has the empty decomposition; a
nonzero indecomposable space is already a factor; otherwise split it into two
nonzero invariant direct summands and apply induction to both smaller dimensions.
Additivity in the lemma shows that every resulting decomposition has exactly
\(\mu(V^{\otimes n})\) null factors. This gives the middle inequality in
\eqref{eq:bounds}; the last follows because each factor has dimension at least
one. The first is the rank bound of the lemma.

Since \(r_s\leq d-1\), division by \(d^n\) gives
\begin{equation}\label{eq:error}
 0\leq1-\frac{b(n)}{d^n}
 \leq2\sum_{s\in S}\left(\frac{r_s}{d}\right)^n
 \leq2t\left(\frac{d-1}{d}\right)^n.
\end{equation}
The right side tends to zero, proving the limit. If \(d=1\), every action is
zero and \(Z=V\) already. If \(d\geq2\), the same geometric decay yields an
exponent satisfying \eqref{eq:threshold}. Its positive lower bound for \(\mu\)
gives a nonzero null summand; any line in that summand has an invariant complement.
\end{proof}

For \(d\geq2\), an explicit sufficient exponent is
\[
 m=1+\left\lfloor\frac{\log(2t)}{\log(d/(d-1))}\right\rfloor.
\]
Indeed \(m\log(d/(d-1))>\log(2t)\), which gives
\(2t(d-1)^m<d^m\). The rank-specific test \eqref{eq:threshold} can be smaller.
Neither algebraic closure nor the absence of a semigroup identity is used.
For a monoid, this statement concerns nonunital representations with singular
identity action; it does not concern ordinary unital modules.

\begin{corollary}\label{cor:length}
Let \(\ell(n)\) be the composition length of \(V^{\otimes n}\), and let
\(c_Z(n)\) be its composition-factor multiplicity of \(Z\). Then
\[
 \mu(V^{\otimes n})\leq c_Z(n)\leq\ell(n)\leq d^n,
 \qquad b(n)\leq\ell(n).
\]
Each of \(\mu(V^{\otimes n})/d^n\), \(c_Z(n)/d^n\), and \(\ell(n)/d^n\)
tends to one, with deficit at most the right side of \eqref{eq:error}.
\end{corollary}
\begin{proof}
Finite-dimensional modules have composition series by induction on dimension,
and each simple factor has positive dimension. A direct summand \(Z^{\oplus\mu}\)
contributes \(\mu\) such factors. Composition length is additive on direct sums,
and every nonzero indecomposable factor has positive length. These facts give all
the inequalities. Apply \eqref{eq:bounds} and divide by \(d^n\).
\end{proof}
The length asymptotic is already Proposition~4.1 of He--Tubbenhauer
\cite[p.~17]{HT25}; the corollary supplies a rank-based proof and error bound.

\section{The related monoid conjecture}
In Conjecture~3A.2, \cite[p.~11]{HT25} requires a tensor power that
``contains a direct summand of the form'' \(\operatorname{Ind}(W)\), where
\(W\ne0\) is a module for the unit group \(G\) and nonunits act as zero.
Here \(\operatorname{Ind}\) has the source's meaning of extension by zero,
not tensor induction. The following gives the nonmodular case without the
right-idempotent-generation assumption of \cite[Theorem~4A.12(c)]{HT26}.

\begin{corollary}\label{cor:monoid}
Let \(M\) be a finite monoid with unit group \(G\), and let \(V\ne0\) be a
finite-dimensional unital \(kM\)-module. Suppose no element other than \(1\)
acts as the identity and \(|G|\) is invertible in \(k\). Put \(S=M\setminus G\),
\(d=\dim V\) and \(r_s=\rank\rho(s)\). For every \(n\geq1\), there is an
invariant decomposition
\[
 V^{\otimes n}=L_n\oplus U_n,\qquad S L_n=0,
 \qquad \dim U_n\leq2\sum_{s\in S}r_s^n.
\]
Consequently Conjecture~3A.2 holds under these hypotheses. If \(b_M(n)\) and
\(b_G(n)\) count indecomposable summands before and after restriction to \(G\),
then
\begin{equation}\label{eq:monoid}
 0\leq b_G(n)-b_M(n)\leq2\sum_{s\in S}r_s^n,
 \qquad b_M(n)/b_G(n)\longrightarrow1.
\end{equation}
When \(G=\{1\}\), over every field,
\(b_M(n)/d^n\to1\); the extra composition-ratio hypothesis of
\cite[Proposition~3C.2, p.~16]{HT25} is unnecessary.
\end{corollary}
\begin{proof}
If \(S\) is empty, take \(L_n=V^{\otimes n}\). Otherwise \(S\) is a two-sided
ideal. If \(\rho(s)\) were invertible, a positive idempotent power \(e\) of
\(s\) would act as the identity. The hypothesis gives \(e=1\), making \(s\)
a unit, a contradiction. Thus every \(r_s<d\).

On \(V^{\otimes n}\), the common kernel \(K\) and image sum \(I\) of the
\(S\)-actions are \(G\)-invariant, because multiplication by units permutes
\(S\). Maschke's theorem \cite[Theorem~4.1.1(i)]{EGH} gives a \(G\)-invariant
complement \(L_n\) of \(K\cap I\) in \(K\), and a \(G\)-invariant
complement \(U_n\) of \(L_n\) containing \(I\). For the latter, first split
\(L_n+I=L_n\oplus I\), then choose a \(G\)-invariant complement of \(L_n+I\)
in the whole space. Every nonunit kills \(L_n\) and maps \(U_n\) into \(I\).
The dimension estimate follows exactly as in Lemma~\ref{lem:split}.
Geometric decay makes \(L_n\ne0\) for all sufficiently large \(n\).

The module \(L_n\) has the same indecomposable summands over \(M\) and \(G\).
Over \(G\), every module is semisimple, so restricting an indecomposable
\(M\)-summand contributes at least one simple \(G\)-summand. Additivity and
Krull--Schmidt \cite[Theorem~3.8.1]{EGH} therefore give
\(0\leq b_G(n)-b_M(n)\leq\dim U_n\).
Each simple \(kG\)-module has dimension at most \(|G|\), since it is a
quotient of the regular module. Hence \(b_G(n)\geq d^n/|G|\); division in
\eqref{eq:monoid} and \(r_s<d\) prove the ratio limit. For trivial \(G\),
\(b_G(n)=d^n\), and the argument applies in every characteristic.
\end{proof}
The leading-asymptotic part of the cross-referenced
\cite[v1, Conjecture~1B.1(1), p.~2; v2, Conjecture~1(1), p.~2]{HTgrowth}
also follows in nonmodular characteristic: the monoid count has the same
asymptotic as its unit-group restriction. Its specified spectral error formulas
are distinct from the rank bound \eqref{eq:monoid}. The modular case with
nontrivial \(G\) is outside this corollary.

\section{A finite-exponent correction}
The conclusion printed in \cite[Proposition~4.4, p.~17]{HT25} is
``\(b(n)=(\dim V)^n\)''. The next example satisfies its uniform
composition-ratio hypothesis but has a nonzero deficit at every positive
exponent. The corrected universal conclusion is \(b(n)\sim d^n\), proved
by Theorem~\ref{thm:main} even without that extra hypothesis.

\begin{corollary}\label{cor:rightzero}
Let \(S=\{x,y\}\) with \(st=t\) for all \(s,t\in S\), and, over any field, set
\[
 \rho(x)=\begin{pmatrix}1&0\\0&0\end{pmatrix},\qquad
 \rho(y)=\begin{pmatrix}1&1\\0&0\end{pmatrix}.
\]
For every \(n\geq1\),
\[
 V^{\otimes n}\cong V\oplus Z^{\oplus(2^n-2)},\qquad b(n)=2^n-1.
\]
Every indecomposable \(kS\)-module other than \(Z\) has null
composition-factor fraction at most \(1/2\). Thus the hypothesis of
Proposition~4.4 holds with \(c=3/4\), while its literal finite-exponent
equality fails. Its asymptotic replacement holds.
\end{corollary}
\begin{proof}
In an arbitrary representation write the two actions as \(P,Q\). The relations
give \(P^2=P\), \(Q^2=Q\), \(PQ=Q\), \(QP=P\); in particular both images
are the same space \(E\). In a decomposition \(W=E\oplus\ker P\),
\[
 P=\begin{pmatrix}1_E&0\\0&0\end{pmatrix},\qquad
 Q=\begin{pmatrix}1_E&C\\0&0\end{pmatrix}.
\]
Separate bases in \(E\) and \(\ker P\) put \(C\) into rank normal form.
It follows that every module is a sum of the displayed two-dimensional module
\(V\), the one-dimensional module \(T\) on which both elements act as one,
and \(Z\). The module \(V\) is indecomposable: its common kernel is zero,
so it has no \(Z\) summand, while \(P\ne Q\) excludes \(T\oplus T\).
Its invariant line \(ke_1\) is \(T\) and its quotient is \(Z\).
Hence \(r(T)=0\) and \(r(V)=1/2\), proving the uniform assertion.

In \(V^{\otimes n}\), both actions have image \(k e_1^{\otimes n}\).
Their defining covectors are \((e_1^*)^{\otimes n}\) and
\((e_1^*+e_2^*)^{\otimes n}\), which are independent. Thus \(\dim K=2^n-2\)
and \(K\cap I=0\). Lemma~\ref{lem:split} gives \(2^n-2\) null summands.
The remaining two-dimensional summand is \(V\), since the restrictions to
\(\langle e_1^{\otimes n},e_2^{\otimes n}\rangle\) are the displayed
matrices. This proves the decomposition and its count.
\end{proof}

\section{Exact examples and coverage}
\begin{corollary}\label{cor:earlierexample}
In \cite[v1, Example~2.2, p.~3; v2, Example~1, p.~3]{HTgrowth}, let
\(S=G_1\cup G_2\cup\{0\}\), with \(G_1=C_2\), \(G_2=S_3\) and
cross-products zero. For the specified complex module \(V=X\oplus Y\),
where \(X\) is the nontrivial \(C_2\)-module and \(Y\) the standard
two-dimensional \(S_3\)-module, every \(n\geq1\) satisfies
\[
 \mu(V^{\otimes n})=3^n-2^n-1,\qquad
 b(n)=3^n-\frac{2^n-(-1)^n}{3}.
\]
\end{corollary}
\begin{proof}
Expand \((X\oplus Y)^{\otimes n}\). All mixed tensor words have zero action,
and their total dimension is \(3^n-2^n-1\). On each pure word the relevant
group identity acts as one, so it has no null summand. The word \(X^{\otimes n}\)
is one-dimensional. For \(Y^{\otimes n}\), the standard \(S_3\)-character on
the identity, transpositions and three-cycles is \((2,0,-1)\); the sum of the
three irreducible characters is \((4,0,1)\). Character inner products give
\(b_{S_3}(n)=(2^{n+1}+(-1)^n)/3\). Adding the pure and mixed counts proves
the formula. These are exact consequences for the source's existing example.
\end{proof}

\begin{proposition}\label{prop:jordan}
For \(d\geq2\), let \(J_d e_i=e_{i-1}\) for \(1\leq i<d\), and
\(J_de_0=0\). Let the semigroup of its positive powers, including zero, act
on \(k^d\). For every \(n\geq1\),
\[
 b(n)=d^n-(d-1)^n,\qquad
 \mu(V^{\otimes n})=d^n-2(d-1)^n+(d-2)^n.
\]
\end{proposition}
These are explicit consequences of classical nilpotent Jordan theory;
the two-factor decomposition is given in
Brualdi \cite[Theorem~4.6(4), p.~47]{Bru85}.
\begin{proof}
In the tensor basis indexed by \(\{0,\ldots,d-1\}^n\), the operator
\(B=J_d^{\otimes n}\) decreases every coordinate by one if all coordinates are
positive, and otherwise gives zero. Nonzero basis images are distinct, so the
basis graph is a disjoint union of directed finite chains. The span of each
chain is invariant under all powers of \(B\). A chain of length \(h\) is
indecomposable: it is generated by its last basis vector, and its commuting
endomorphisms form \(k[x]/(x^h)\); this ring has only the idempotents zero and
one, since its nonunits form the nilpotent ideal \((x)\). A nontrivial invariant
direct-sum decomposition would give another commuting idempotent.

The number of chains is \(\dim V^{\otimes n}-\rank B=d^n-(d-1)^n\).
A chain of length \(h\) contributes
\(h-2\max(h-1,0)+\max(h-2,0)\) to
\(\dim V^{\otimes n}-2\rank B+\rank B^2\), which is one for \(h=1\) and
zero otherwise. Since \(\rank B^2=(d-2)^n\), this counts the null chains and
proves the second formula.
More generally, the number of blocks of size exactly \(j\), for \(1\leq j\leq d\), is
\begin{equation}\label{eq:jordanall}
 m_j(n)=(d-j+1)^n-2(d-j)^n+\max(d-j-1,0)^n.
\end{equation}
Indeed \(\rank B^q=\max(d-q,0)^n\). A block of size \(h\) contributes
\(\max(h-j+1,0)-2\max(h-j,0)+\max(h-j-1,0)\) to the corresponding
second difference of ranks, equal to one exactly when \(h=j\).
\end{proof}
In particular, the first tensor power need not contain \(Z\), whereas the second
power in this family has exactly two copies of \(Z\). For the operator family
\(\{0,J_2\}\) on \(k^2\), \(\mu=0\) and the sum of the ranks is one.
Thus the constant two in the universal estimate
\(\mu(W)\geq\dim W-2\sum_s\rank A_s\) cannot be decreased.

\begin{corollary}\label{cor:fusion}
The matrix in Example~4.3 and Figure~9 of \cite[p.~17]{HT25} yields, for \(n\geq2\),
\[
 b(n)=6^n-7,\qquad [V^{\otimes n}:Z]_b=6^n-8.
\]
\end{corollary}
\begin{proof}
Use column vectors and the ordered
dimensions \((1,6,1,8)\), with the third vertex equal to \(Z\). The printed
matrix is
\[
 A=\begin{pmatrix}0&0&0&0\\1&0&0&0\\0&28&6&40\\0&1&0&1\end{pmatrix}.
\]
Starting from \(v_0=(1,0,0,0)^T\), put \(v_{n+1}=Av_n\). Direct multiplication
gives \(v_2=(0,0,28,1)^T\). If
\(v_n=(0,0,6^n-8,1)^T\), multiplication gives
\(v_{n+1}=(0,0,6(6^n-8)+40,1)^T=(0,0,6^{n+1}-8,1)^T\).
The sum of the coordinates is \(6^n-7\), and the third coordinate is
\(6^n-8\), as claimed. The calculation concerns the published matrix.
\end{proof}

\begin{center}\small
\begin{tabularx}{\textwidth}{@{}>{\raggedright\arraybackslash}p{.31\textwidth}>{\raggedright\arraybackslash}X@{}}
\toprule Source item in \cite{HT25} & Coverage \\ \midrule
Conjecture 4.2, p.~17 (\texttt{ht2508}) & Proved by Theorem~\ref{thm:main}, including the ratio limit. \\
Right-idempotent-generated case, \cite[Theorem 4A.14(e)--(f)]{HT26} & Reproved; prior proof He--Tubbenhauer, 19 December 2025. Extended to all finite singular actions over any field. \\
\S1(1)--(4), pp.~1--2 & In the singular-semigroup setting: exponential error bound, absolute deficit, and null multiplicity from \eqref{eq:bounds}--\eqref{eq:error} and Corollary~\ref{cor:length}. \\
Proposition 4.1, p.~17 & Reproved; prior proof He--Tubbenhauer (2025). Rank error bound added; its \(\mathcal L\)-class bound is not covered. \\
Example 4.3 / Figure 9, p.~17 & Prior matrix He--Tubbenhauer (2025); exact recurrence proved in Corollary~\ref{cor:fusion}. \\
\cite[v1, Example 2.2; v2, Example 1, p.~3]{HTgrowth} & Prior example and asymptotic He--Tubbenhauer (2025); exact null and total counts proved in Corollary~\ref{cor:earlierexample}. \\
Proposition 4.4, p.~17, literal finite-\(n\) equality & Refuted-and-corrected by Corollary~\ref{cor:rightzero}: \(b(n)=2^n-1\) meets its hypotheses; the replacement \(b(n)\sim d^n\) holds without the ratio hypothesis. \\
Conjecture 3A.2, p.~11 & Proved when \(|G|\) is invertible in \(k\), Corollary~\ref{cor:monoid}; modular nontrivial \(G\) not covered. \\
Proposition 3C.2, p.~16; Proposition 3C.1 & Reproved; prior proof He--Tubbenhauer (2025); the diagram case over all characteristics is already \cite[Example~4A.13(b)]{HT26}. The general trivial-unit-group case needs no composition-ratio hypothesis. \\
\cite[v1, Conjecture 1B.1(1); v2, Conjecture 1(1), p.~2]{HTgrowth} & Leading-asymptotic assertion proved in nonmodular characteristic; its specified spectral rates and general modular case not covered. \\
\cite[v1, \S4C, p.~8; v2, \S4.3]{HTgrowth} & Its noninjectivity examples remain valid. The natural \(\{1,x,0\}\) module with \(x=J_2\) is covered by Proposition~\ref{prop:jordan} and Corollary~\ref{cor:monoid}; the stated characteristic-31 \(T_3\times N\) example is covered by the monoid corollary, not an exact decomposition formula. \\
\S2 exact character formulas; \S3 subgroup-specific multiplicities; OEIS A009766, A379838, A000085 & Not covered. The monoid corollary compares total counts, not those individual formulas. \\
General multiplicity identity; nilpotent blocks & Reproved; prior principle Juteau--Mautner--Williamson (2009), semigroup criterion McAlister (1971), two-factor blocks Brualdi (1985). Equation~\eqref{eq:jordanall} is an explicit rank-difference consequence. \\
\bottomrule
\end{tabularx}
\end{center}

\subsection*{Statement and verification coverage}
The pinned Lean declaration \texttt{result} in
\href{https://github.com/the-omega-institute/trureturing/blob/15070fa5f71c4237050acc7ebe058958cbc04e6e/D5/S3/HomologicalAlgebra/TensorNullSummand.lean}{\texttt{TensorNullSummand.lean}} \cite{Lean} proves
\texttt{claim}: an algebraically closed field and a nonempty finite semigroup
without identity, with every action singular, imply a positive null
tensor summand, existence of a summand-count sequence, and the ratio limit for
every such sequence. Its definitions \texttt{HasNullSummand},
\texttt{IsSummandCount}, and \texttt{tensorAction} specify an actual invariant
complement, internal indecomposable decompositions, and diagonal tensor actions.
In that same file, \texttt{exists\_null\_tensor\_power},
\texttt{exists\_summand\_count} and \texttt{tensor\_counts\_tendsto} prove the
existence and count-limit statements over any field; their hypotheses do not
exclude a semigroup identity. The tensor lemma \texttt{tensor\_power\_range\_le}
gives a rank upper bound; the exact tensor-rank identity has the written proof
in Theorem~\ref{thm:main}. The maximal-multiplicity formula, monoid corollary,
right-zero decomposition, earlier semigroup example, Jordan formulas and fusion recurrence have the written
proofs above. The accompanying \texttt{verify.py}
checks all 1023 nonempty families of singular binary \(2\times2\) matrices at
\(n=1,2,3\); all null-image projections at \(n=1\); the full family's bound at
\(n=5\); 25 Jordan cases with \(2\leq d\leq6\), \(1\leq n\leq5\); and the
published matrix through \(n=12\). It also checks every block size in those
25 Jordan cases and explicit right-zero decompositions for \(1\leq n\leq7\).
The earlier \(C_2\cup S_3\cup\{0\}\) example is checked through \(n=12\).
These finite checks have exactly those ranges.

\begin{thebibliography}{99}
\raggedright
\small
\bibitem{HT25} D. He and D. Tubbenhauer,
\emph{Tensor powers of representations of (diagram) monoids},
arXiv:2508.04054v1 (6 August 2025), especially pp.~1, 11, 17.
\url{https://arxiv.org/abs/2508.04054v1}.
\bibitem{HT26} D. He and D. Tubbenhauer,
\emph{Affine diagram categories, algebras and monoids},
arXiv:2512.05510v2 (19 December 2025), Theorem~4A.14, p.~20;
revised v3 (20 October 2026), Theorems~4A.12 and 4A.14, pp.~21--22.
\url{https://arxiv.org/abs/2512.05510v2};
\url{https://arxiv.org/abs/2512.05510v3}.
Forthcoming in \emph{Journal of Algebra},
\url{https://doi.org/10.1016/j.jalgebra.2026.09.034}.
\bibitem{HTgrowth} D. He and D. Tubbenhauer,
\emph{Growth problems for representations of finite monoids},
arXiv:2502.02849v1 (5 February 2025), Conjecture~1B.1 and Theorem~3A.2;
v2 (8 November 2025), Conjecture~1 and Theorem~1.
\url{https://arxiv.org/abs/2502.02849v1};
\url{https://arxiv.org/abs/2502.02849v2}.
\bibitem{JMW} D. Juteau, C. Mautner and G. Williamson,
\emph{Parity sheaves}, arXiv:0906.2994v1 (16 June 2009),
\S3.2.1, p.~23, (3.1); v3 (17 February 2014), Lemma~3.1, pp.~22--23.
\emph{J. Amer. Math. Soc.} \textbf{27} (2014), 1169--1212.
\url{https://arxiv.org/abs/0906.2994v1};
\url{https://arxiv.org/abs/0906.2994v3};
\url{https://doi.org/10.1090/S0894-0347-2014-00804-3}.
\bibitem{McA71} D. B. McAlister,
\emph{The category of representations of a completely 0-simple semigroup},
\emph{J. Austral. Math. Soc.} \textbf{12} (1971), 193--210,
Definition~1.4 and Corollary~2.6.
\url{https://doi.org/10.1017/S1446788700009460}.
\bibitem{Bru85} R. A. Brualdi,
\emph{Combinatorial verification of the elementary divisors of tensor products},
\emph{Linear Algebra Appl.} \textbf{71} (1985), 31--47,
Lemma~3.3 and Theorem~4.6(4).
\url{https://doi.org/10.1016/0024-3795(85)90233-2}.
\bibitem{COT} K. Coulembier, V. Ostrik and D. Tubbenhauer,
\emph{Growth rates of the number of indecomposable summands in tensor powers},
\emph{Algebr. Represent. Theory} \textbf{27} (2024), 1033--1062,
Theorem~1.4(a), Remark~1.5(a).
\url{https://arxiv.org/abs/2301.00885v2};
\url{https://doi.org/10.1007/s10468-023-10245-7}.
\bibitem{LZ} M. Lu and Y. Zhang,
\emph{Growth rates of indecomposable summands in tensor powers of
representations of quivers}, arXiv:2605.03506v1 (5 May 2026), Theorem~A.
\url{https://arxiv.org/abs/2605.03506v1}.
\bibitem{EGH} P. Etingof, O. Golberg, S. Hensel, T. Liu, A. Schwendner,
D. Vaintrob and E. Yudovina, \emph{Introduction to representation theory},
\emph{Student Mathematical Library} \textbf{59}, American Mathematical Society,
2011, Theorems~3.8.1 and 4.1.1(i).
\url{https://math.mit.edu/~etingof/reprbook.pdf}.
\bibitem{Lean} \emph{TensorNullSummand.lean},
declaration \texttt{result} in namespace
\nolinkurl{D5.S3.HomologicalAlgebra.TensorNullSummand},
revision \texttt{15070fa5f71c4237050acc7ebe058958cbc04e6e}.
\url{https://github.com/the-omega-institute/trureturing/blob/15070fa5f71c4237050acc7ebe058958cbc04e6e/D5/S3/HomologicalAlgebra/TensorNullSummand.lean}.
\end{thebibliography}
\noindent\small Licensed under Creative Commons Attribution 4.0 International
(CC BY 4.0): \url{https://creativecommons.org/licenses/by/4.0/}.
\end{document}
